% engines: lualatex
% expect-warnings: 5
% expect-errors: 0
% passes: 2
%
% animate needs two runs.
\documentclass{article}
\usepackage[margin=1cm]{geometry}
\usepackage{animate}
\usepackage[kernel = lua]{../../latex/pgf-go}
\newcommand\case[1]{\par\medskip\noindent\textbf{#1}\par\nopagebreak}
\pagestyle{empty}

% With layers, every stone is drawn once, in a layer that animate shows while
% the stone does not change. \layers{<options>}{<file>}{<frames>} draws each
% frame twice: whole (as with layers = false) and made of its layers, one on
% top of the other. Both must look the same.
\ExplSyntaxOn
\NewDocumentCommand \layers {mmm} {
	\group_begin:
		\pgfgo_sgfanimation_setup:nn {#1} {#2}
		\pgfgo_sgfanimation_layers:
		\clist_map_inline:nn {#3} {
			\pgfgo_sgfanimation_frame:n {##1}\enspace
			\pgfgo_sgfanimation_layers_frame:n {##1}\quad
		}
	\group_end:
}

% Every layer must have the size of the board: \samesize{<options>}{<file>}.
\msg_new:nnn {sgf-animation-layers-test} {size}
	{Layer~#1~is~#2,~but~layer~0~is~#3.}
\tl_new:N \l__test_size_tl
\tl_new:N \l__test_first_size_tl
\NewDocumentCommand \samesize {mm} {
	\group_begin:
		\pgfgo_sgfanimation_setup:nn {#1} {#2}
		\pgfgo_sgfanimation_layers:
		\int_step_inline:nnn {0} {\g__pgfgo_sgfanim_layer_int - 1} {
			\tl_set:Ne \l__test_size_tl {
				\dim_use:N \box_wd:c {g__pgfgo_sgfanim_##1_box} ~ x ~
				\dim_eval:n {\box_ht:c {g__pgfgo_sgfanim_##1_box} + \box_dp:c {g__pgfgo_sgfanim_##1_box}}
			}
			\int_compare:nNnTF {##1} = {0}
				{\tl_set_eq:NN \l__test_first_size_tl \l__test_size_tl}
				{
					\tl_if_eq:NNF \l__test_size_tl \l__test_first_size_tl
						{\msg_error:nneee {sgf-animation-layers-test} {size} {##1} {\l__test_size_tl} {\l__test_first_size_tl}}
				}
		}
	\group_end:
}

% The layers of each stone: \stonelayers{<options>}{<file>} writes
% <layer>:<first frame>x<frames>.
\NewDocumentCommand \stonelayers {mm} {
	\group_begin:
		\pgfgo_sgfanimation_setup:nn {#1} {#2}
		\pgfgo_sgfanimation_layers:
		\seq_map_inline:Nn \l__pgfgo_sgfanim_stones_seq {\__test_stone:nnn ##1}
	\group_end:
}
\cs_new:Npn \__test_stone:nnn #1#2#3 {#1:#2x#3~}
\ExplSyntaxOff

\begin{document}

% game.sgf: 9x9, black C7, C6, C5, white G3; then B E5, W G7, B pass,
% W C3 (comment "Main line."), B G5.

\case{The animation by layers (pdftoppm only draws its controls)}
\sgfanimation[comments, goban = {captures}]{game.sgf}

\case{game.sgf, whole and by layers: frames 0, 3 (a pass) and 5; numbers: frame 5}
\layers{}{game.sgf}{0, 3, 5}

\layers{numbers}{game.sgf}{5}

\case{Captures (frames 6 and 7 of the manual's game: 7 captures 2), with comments}
\layers{comments, goban = {captures}}{../../doc/en/game.sgf}{6, 7}

\case{setup.sgf: AB, AW and AE after move 3; with captures and without (AE needs captures: frame 4 keeps the stone)}
\layers{goban = {captures}}{setup.sgf}{3, 4}

\layers{}{setup.sgf}{4}

\case{variations.sgf: marks of the file on stones and on empty points (frame 2); variation = 2 (frame 3)}
\layers{}{variations.sgf}{2}

\layers{variation = 2}{variations.sgf}{3}

\case{ko.sgf with captures = fade and captured marks = remove: frames 2 and 3}
\layers{goban = {captures = fade, captured marks = remove}}{ko.sgf}{2, 3}

\case{Layers of the stones of game.sgf: the setup stones (1 to 4) last 6 frames}
\stonelayers{}{game.sgf}

\case{Layers of ko.sgf with captures: C3 is played, captured and played again}
\stonelayers{goban = {captures}}{ko.sgf}

\samesize{comments, goban = {captures}}{../../doc/en/game.sgf}
\samesize{comments}{setup.sgf}
\samesize{}{variations.sgf}

\end{document}
